\documentclass{article}
\usepackage[utf8]{inputenc}
\usepackage[T1]{fontenc}
\usepackage[margin=2cm]{geometry}
\usepackage{hyperref}
\sloppy
\begin{document}
\section*{Historical source}
Encyclopaedia Britannica scan, source row 2493. Parse 3.5 extracted the text. Super 3.5 VL corrected it from an enlarged crop of each block. It can still contain recognition errors.
Page headers, footers, and page numbers are removed.

\subsection*{CHEMISTRY (E0247)}
CHEMISTRY.

Magnella,
Sc.
Magnesia,
\&c.

Preparation.

1353
Properties

1354
Action of
air, \&c.

into sulphate of magnesia. It is soluble in 20 parts of cold water. Boiling water dissolves a greater proportion, and from this it crystallizes on cooling. Exposed to heat, this salt becomes viscid, and by calcination it loses 0.45 of its weight. If the heat be increased, it is decomposed; the acid is driven off, and the pure earth remains behind.

The component parts of this fall are,

Preparation.

1356
Properties.

1357
Action of
heat.

1358
History and preparation.

1359
Properties.

3. It is deliquescent in the air, and is soluble in its own weight of cold water. It is more soluble in boiling water, in which it crystallizes on cooling, but it

1356
Of heat.

1362
Composition.

3. Sulphite of Magnesia.

1. The compound of sulphurous acid and magnesia is formed by passing sulphurous acid gas into two parts of water, with one of carbonate of magnesia. A violent effervescence takes place, with the evolution of heat. The sulphite of magnesia is formed, and precipitated to the bottom in the state of powder; but with an excess of acid it is re-dissolved, and crystallizes.

2. The crystals of sulphite of magnesia are in the form of deprived transparent tetrahedrons. It has a mild earthy taste, which soon becomes feebly sulphu-

3. It effloresces in the air, and is slowly converted into sulphate of magnesia. It is soluble in 20 parts of cold water. Boiling water dissolves a greater proportion, and from this it crystallizes on cooling. Exposed to heat, this salt becomes viscid, and by calcination it loses 0.45 of its weight. If the heat be increased, it is decomposed; the acid is driven off, and the pure earth remains behind.

The component parts of this fall are,

r. This triple salt is formed by decomposing the sulphite of ammonia by magnesia, or the sulphite of magnesia by ammonia, in solution in the cold; or, by mixing together the solutions of the two salts.

2. This salt is in transparent crystals, the form of which has not been determined. When it is exposed to the air, it is converted into sulphate of ammonia and magnesia. It is less soluble in water than either of the two sulphites of which it is formed. By the action of heat, sulphurous acid is given out, acidulous sulphite of ammonia is sublimed, and there remains behind pure

5. Nitrate of Magnesia.

1358
History and preparation.

1. This compound of nitric acid and magnesia was formerly called nitre with base of magnesia, and magnesian saltpetre. It is formed by the direct combination of the acid with the earth. By evaporation it is crystallized.

2. This fatty crytallizes in four-fied rhomboidal prisms, whose fummits are oblique or truncated. Sometimes it is in the form of small needles combined in groups. The taste is penetrating and bitter. The specific gravity is 1.736.

3. It is deliquescent in the air, and is soluble in its own weight of cold water. It is more soluble in boiling water, in which it crystallizes on cooling; but it can only be obtained in regular crystals by slow evaporation from its solution in cold water.

4. By the action of heat it undergoes the watery fusion; the water is driven off, and it becomes dry. It is decomposed in a strong heat, gives out a little oxygen gas, then nitrous gas, and at last the nitric acid. The pure earth remains behind.

The component parts of this fault are, according to

\textbackslash{}begin\{tabular\}\{cccc\}
\textbackslash{}thead\{\textbackslash{}mbox\{\}\} \& \textbackslash{}thead\{Bergman.\} \& \textbackslash{}thead\{Kirwan.\} \& \textbackslash{}thead\{Magne\}\textbackslash{}\textbackslash{}
Acid \& 43 \& 46 \& \textbackslash{}mbox\{\}\textbackslash{}\textbackslash{}
Magnella \& 27 \& 22 \& \textbackslash{}mbox\{\}\textbackslash{}\textbackslash{}
Water \& 30 \& 32 \& \textbackslash{}mbox\{\}\textbackslash{}\textbackslash{}
100 \& 100 \& 100 \& \textbackslash{}\textbackslash{}
\textbackslash{}end\{tabular\}

6. Nitrate of Ammonia and Magnesia.

1. This triple salt is formed, either by the direct combination of the solutions of nitrate of ammonia, and nitrate of magnesia, by which the salt is obtained pure and crystallized; or, by partially decomposing the nitrate of ammonia by magnesia, or the nitrate of magnesia by ammonia.

2. The cryftals of this salt are in the form of fine prifms. It has a bitter, acrid, and ammoniacal taste. It is less deliquescent in the air than either of the confluent falts, and less soluble in water. It requires 11 parts of cold water to difolve it, but less of boiling water. It cryftallizes on cooling.

When it is rapidly heated, it inflames spontaneously. When slowly heated in close vessels, it gives out oxygen gas, azotic gas, a greater proportion of water than it contains, nitrous gas, and nitric acid, without the faintest trace of ammonia; which shows that it is decomposed, that the hydrogen combines with the oxygen of the acid, and forms water.

The component parts of this fall are,

100 *

7. Nitrite of Magnesia.

Nothing is known of the properties of this fault.

8. Muriate of Magnesia.

This compound of muriatic acid and magnesia found was formerly called marine salt of magnesia, and was not confounded with the muriate of lime, with which it is frequently accompanied. The difference between the two salts was first pointed out by Dr Black, and Bergman afterwards examined the nature and properties of muriate of magnesia. The salt is obtained by dissolving magnesia in muriatic acid till they are saturated, and then evaporating the solution. Small irregular crystals are obtained. This salt exists in the waters of the ocean, and in mineral waters, along with the muriates of soda and lime.

2. It is extremely difficult to obtain the muriate of proper magnetic in any regular form. It is either in the state of powder, or very small regular needles, or in a kind of jelly. It has a disagreeable bitter taste. The specific gravity is 1.601.

3. It is very deliquescent in the air. Cold water readily dissolves its own weight, and it is still more soluble in boiling water.

4. It is completely decomposed by heat ; the acid is driven off, and the pure earth remains behind *.

Acid 34 34-59
Magnesia 41 31.07
Water 25 34.38

\textbackslash{}begin\{tabular\}\{cc\}
Acid \& 34-19\textbackslash{}\textbackslash{}
Magnella \& 41\textbackslash{}\textbackslash{}
Water \& 25\textbackslash{}\textbackslash{}
\textbackslash{}end\{tabular\}


\section*{Source record}
\begin{description}
\item[Scan ID] \texttt{britannica:309f765b1ecc83e0a7f7762e44b4a927a7555443:train:2493}
\item[Image file] \texttt{page-2493-84519542efce.jpg}
\item[Image SHA-256] \texttt{84519542efceeb7511318fcd5cf2213ed3dc634dd20638613d3ba33a500c2393}
\item[Dataset] NationalLibraryOfScotland/encyclopaedia\_britannica\_illustrated at revision \texttt{309f765b1ecc83e0a7f7762e44b4a927a7555443} (CC0-1.0)
\item[Extraction] model \texttt{nvidia/nemotron-3.5-parse}, raw reply \texttt{extract/raw/page-2493-84519542efce.json}
\item[Raw reply SHA-256] \texttt{0ba2d70cd1f7c113f8e899cc3ddb1b05e6ebed12cf9c6d145301c4e6998906d4}
\item[Correction] model \texttt{nvidia/Nemotron-3.5-Super-VL}, block crops, raw replies \texttt{correct/raw/crop\_nothink/page-2493-84519542efce-b*.json}; parse: 4, super\_vl: 43
\item[Running header] none
\end{description}
\end{document}
